\emph{Preliminaries.} The blades are a basis (Theorem \claimref{synthesis.signed-directions-and-the-integer-clifford-algebra}); the adjoint and the scalar part are those of Theorem \claimref{synthesis.chsh-bound-for-two-subsystems-of-a-clifford-algebra-by-a-rational-certificate}, with $\langle xy\rangle_{0} = \langle yx\rangle_{0}$. By Proposition \claimref{tower.ladder-pairs-satisfy-car-vectors}, $v_{j}^{2} = w_{j}^{2} = 0$, $v_{j}w_{j} + w_{j}v_{j} = 1$, and ladder elements of different modes anticommute; $w_{j} = v_{j}^{\dagger}$ since $\dagger$ fixes $e_{j}$ and negates $f_{j}$.

\emph{(i).} $w_{j}v_{j}$ is a self-adjoint idempotent ($w_{j}v_{j}w_{j}v_{j} = w_{j}(1 - w_{j}v_{j})v_{j} = w_{j}v_{j}$), the two commute, so $P_{0}$ is a self-adjoint idempotent with $w_{j}P_{0} = 0$ and $P_{0}v_{j} = 0$. The relations $s_{i}^{\dagger}s_{j} = \delta_{ij}P_{0}$ and $\sum_{i}s_{i}s_{i}^{\dagger} = 1$ follow from these and are verified by the check. Hence $y_{r}^{\dagger}y_{r'} = (r \cdot r')P_{0}$ and $y_{r}P_{0} = y_{r}$, so $\Pi_{r}\Pi_{r'} = (r \cdot r')\,y_{r}y_{r'}^{\dagger}/((r \cdot r)(r' \cdot r'))$: $\Pi_{r}$ is a self-adjoint idempotent, and $\Pi_{r}\Pi_{r'} = 0$ for orthogonal vectors. $P_{0}xP_{0}$ is a multiple of $P_{0}$ for every $x$ (verified on the blades, with $\langle P_{0}\rangle_{0} = \tfrac14$), so $\Pi_{r}x\Pi_{r}$ is a multiple of $y_{r}P_{0}y_{r}^{\dagger}$, that is of $\Pi_{r}$, and taking scalar parts with $\langle\Pi_{r}\rangle_{0} = \langle y_{r}^{\dagger}y_{r}\rangle_{0}/(r \cdot r) = \tfrac14$ gives the factor $4\langle\Pi_{r}x\rangle_{0}$. For four pairwise orthogonal vectors the sum $S$ of the four atoms satisfies $Sy_{r} = y_{r}$ for each of them; the four vectors are a basis of $\Q^{4}$, so $Ss_{i} = s_{i}$ and $S = S\sum_{i}s_{i}s_{i}^{\dagger} = 1$. The nine contexts are the rows of
\[
\begin{array}{llll}
(0,0,0,1) & (0,0,1,0) & (1,1,0,0) & (1,-1,0,0) \\
(0,0,0,1) & (0,1,0,0) & (1,0,1,0) & (1,0,-1,0) \\
(1,-1,1,-1) & (1,-1,-1,1) & (1,1,0,0) & (0,0,1,1) \\
(1,-1,1,-1) & (1,1,1,1) & (1,0,-1,0) & (0,1,0,-1) \\
(0,0,1,0) & (0,1,0,0) & (1,0,0,1) & (1,0,0,-1) \\
(1,-1,-1,1) & (1,1,1,1) & (1,0,0,-1) & (0,1,-1,0) \\
(1,1,-1,1) & (1,1,1,-1) & (1,-1,0,0) & (0,0,1,1) \\
(1,1,-1,1) & (-1,1,1,1) & (1,0,1,0) & (0,1,0,-1) \\
(1,1,1,-1) & (-1,1,1,1) & (1,0,0,1) & (0,1,-1,0)
\end{array}
\]
Each row is an orthogonal set, there are eighteen distinct vectors, and each occurs in exactly two rows (verified by the check, together with the idempotent, orthogonality and sum relations for the eighteen atoms themselves). If $g$ took values in $\{0, 1\}$ with sum $1$ on each row, the sum of $g$ over all nine rows would be $9$; it is also twice the sum of $g$ over the eighteen atoms, an even number. The check confirms this by trying all $2^{18}$ assignments. A function on all atoms summing to $1$ on every context would restrict to such a $g$.

\emph{(ii).} $X_{1}, Z_{1}, X_{2}, Z_{2}$ are self-adjoint blades up to sign of square $1$ with $X_{j}Z_{j} = -Z_{j}X_{j}$ and each of $X_{1}, Z_{1}$ commuting with each of $X_{2}, Z_{2}$ (Theorem \claimref{tower.commuting-copies} and clause (vii) of the CHSH theorem). Products of commuting self-adjoint elements of square $1$ are self-adjoint of square $1$; the entries of each row and each column commute pairwise (for the third row and column this uses two anticommutations). Row products: $X_{1}X_{2}(X_{1}X_{2}) = 1$, $Z_{2}Z_{1}(Z_{1}Z_{2}) = 1$, $(X_{1}Z_{2})(Z_{1}X_{2})(X_{1}X_{2}Z_{1}Z_{2}) = -1$; column products: $X_{1}Z_{2}(X_{1}Z_{2}) = 1$, $X_{2}Z_{1}(Z_{1}X_{2}) = 1$, $(X_{1}X_{2})(Z_{1}Z_{2})(X_{1}X_{2}Z_{1}Z_{2}) = 1$; all six are computed by the check. Multiplying the six conditions on $\nu$, each $\nu(a)$ occurs twice, so the left side is $1$, and the right side is $-1$.

\emph{(iii).} $\Gamma = p_{1}p_{2}$ is self-adjoint of square $1$ and commutes with every even element, so $Q_{\pm} = \tfrac12(1 \pm \Gamma)$ are self-adjoint idempotents central in $\mathcal{A}^{+}$. $s_{1} = p_{1}$ and $s_{2} = e_{1}f_{2}$ are even self-adjoint blades of square $1$ that anticommute, and $s_{1}s_{2}$ is anti-self-adjoint. The corner $Q_{+}\mathcal{A}^{+}$ has dimension $4$ and basis $Q_{+}, Q_{+}s_{1}, Q_{+}s_{2}, Q_{+}s_{1}s_{2}$ (computed by the check). Its self-adjoint elements are $c = Q_{+}(t + xs_{1} + zs_{2})$, with $c^{2} = Q_{+}(t^{2} + x^{2} + z^{2} + 2t(xs_{1} + zs_{2}))$; so its self-adjoint idempotents are $0$, $Q_{+}$ and the elements $c(x, z) = Q_{+}(\tfrac12 + xs_{1} + zs_{2})$ with $x^{2} + z^{2} = \tfrac14$. Define $h(0) = 0$, $h(Q_{+}) = 1$, and $h(c(x, z)) = 1$ if $z > 0$, or $z = 0$ and $x > 0$, and $0$ otherwise; then $h(c(x, z)) + h(c(-x, -z)) = 1$. For a self-adjoint idempotent $E$ of $\mathcal{A}^{+}$ the element $Q_{+}E$ is a self-adjoint idempotent of the corner; put $g(E) = h(Q_{+}E)$, so $g(1) = 1$. If $EE' = 0$ then $Q_{+}E$ and $Q_{+}E'$ are self-adjoint idempotents of the corner with product $0$ and their sum is one too. If one of them is $0$ additivity is clear. If one is $Q_{+}$ the other is $0$. If both are of the form $c(x, z)$, $c(x', z')$, the sum has $t = 1$, so it is $Q_{+}$ and $(x', z') = (-x, -z)$: the values of $h$ add to $1 = h(Q_{+})$. Finally the atom $\Pi_{(1,1,1,1)}$ has a nonzero coefficient on the odd blade $e_{1}$ (computed), so the eighteen atoms are not all even.
